\section{Homología celular}
\label{subsec:homcelular}
 Recordamos que unos de los espacios que tienen especial interés para nosotros, junto con las variedades diferenciables, son los CW-complejos. 
 
 \iffalse
 Como hemos visto, toda variedad diferenciable admite una estructura de CW-complejo vía una función de Morse. Luego la homología celular nos permite computar los grupos de homología de cualquier variedad diferenciable siempre que tengamos una función de Morse. {\color{blue} pero el problema aquí es que na funcion de Morse no nos dice como se pegan no?}
 \fi

 Un espacio triangulado constituye un CW-complejo, al fin y al cabo, cada uno de los $C^k_{\alpha}$ que se unen para construir una triangulación es una $k$-celda. Es de esperar entonces que la \textit{homología celular} se construya de manera similar a la simplicial, aunque ahora en vez de $k$-poliedros tendremos pegados a un espacio $X$ de manera adecuada tendremos $k$-celdas pegadas entre si por las funciones de pegado.
 
Vamos a introducir primero algunos conceptos para evitar posibles lagunas técnicas en el paso de triangulaciones a CW-comlejos en general.
 \iffalse
 \begin{definition}
    Un \textit{complejo de cadena} es una familia de pares $(A_k, \partial_k: A_k\to A_{k-1})$ de grupos abelianos y homomorfismos de grupos llamados \textit{operadores de borde} tales que $\partial_{k-1}\circ \partial_k=0$ para todo $k$.
 \end{definition}
\fi

 \begin{definition}
    Dados dos complejos de cadena $(A_k, \partial_k)$ y $(B_k, \partial_k')$, llamamos \textit{morfismo de cadena} a una familia $f_k$ de homomorfismos tales que $\partial_k'\circ f_k = f_{k-1}\circ \partial_k$ para todo $k$.
 \end{definition}

 Esto resulta en el diagrama siguiente
 \[
 \begin{array}{ccccccccc}
    \dots & \overset{\partial_{k+2}}{\longrightarrow} & A_{k+1} & \overset{\partial_{k+1}}{\longrightarrow} & A_{k} & \overset{\partial_{k}}{\longrightarrow} & A_{k-1} & \overset{\partial_{k-1}}{\longrightarrow} & \dots \\
    & & f_{k+1} \huge\downarrow  ~~~ & & f_k \huge\downarrow ~~ & & f_{k-1} \huge\downarrow ~~~ & & \\
    \dots & \overset{\partial_{k+2}}{\longrightarrow} & B_{k+1} & \overset{\partial_{k+1}}{\longrightarrow} & B_{k} & \overset{\partial_{k}}{\longrightarrow} & B_{k-1} & \overset{\partial_{k-1}}{\longrightarrow} & \dots
 \end{array}
\]

\begin{definition}
    Sea $f \colon (A_\bullet, \partial) \to (B_\bullet, \partial')$ un morfismo de cadena. Como $f \circ \partial = \partial' \circ f$, se tiene que
    $f(\ker \partial_k) \subset \ker \partial'_k$ y $f(\operatorname{im} \partial_{k+1}) \subset \operatorname{im} \partial'_{k+1}$. Por tanto, \( f \) desciende al cociente, dando lugar a un homomorfismo $f_* \colon H_k(A_\bullet, \partial) \to H_k(B_\bullet, \partial')$ que llamamos \textit{homomorfismo inducido en homología}.

\end{definition}
\begin{remark}
    Es claro que $(g \circ f)_* = g_* \circ f_*$,
    para morfismos $(A_\bullet, \partial) \xrightarrow{f} (B_\bullet, \partial') \xrightarrow{g} (C_\bullet, \partial'')$.
\end{remark}

\bigskip

\begin{definition}
    Sea $X$ un CW-complejo. Orientamos cada una de la $k$-celdas de $X$ y definimos
    \begin{equation*}
        C^{CW}_k(X) = \oplus_{\alpha\in\Lambda_k} \mathbb{Z} \langle D^k_{\alpha}\rangle,
    \end{equation*}
    y los operadores de borde $\partial_k : C^{CW}_k(X) \to C^{CW}_{k-1}(X)$ vienen dados por 
    \begin{equation*}
        \partial_k(D^k_{\alpha}) = \sum_{\beta\in\Lambda_{k-1}}[D^k_{\alpha}:D^{k-1}_{\beta}]D^{k-1}_{\beta}
    \end{equation*}
\end{definition}


Igual que en homología simplicial, los elementos de $C^{CW}_k(X)$ son sumas formales de $k$-celdas con coeficientes enteros. Ahora, los operadores de borde siguen llevando una $k$-celda a su borde, solo que esta vez dicho borde será (teniendo en cuenta la orientación) la imagen de su función de pegado. 

Seamos más claros, sean $i_{\alpha}(D^k_{\alpha})$ una $k$-celda de un CW-complejo $X$ y $\varphi_{\alpha}: \partial D_{\alpha}^k \cong \mathbb{S}^{k-1} \to X_{k-1}$ su mapa de pegado. Entonces, $[D^k_{\alpha}:D^{k-1}_{\beta}]$ mide cuántas veces (pueden ser cero) $\varphi_{\alpha}(\partial D_{\alpha}^k)$ recorre completamente (es un entero) la $(k-1)$-celda $i_{\beta}(D^{k-1}_{\beta})$ y en qué sentido, si este coincide con la orientación de $i_{\beta}(D^{k-1}_{\beta})$ el signo será positivo y si difiere negativo. Así, obtenemos una suma formal de elementos de $(k-1)$-celdas, es decir, un elemento de $C^{CW}_{k-1}(X)$. 

Más formalmente, el operador de borde $\partial_k$ realiza el siguiente recorrido. La función de pegado induce un homomorfismo en homología 
\[
(\varphi_{\alpha})_* : H_{k-1}(\partial D_{\alpha}^k)\to H_{k-1}(X_{k-1}).
\]
Este homomorfismo lleva la \textit{clase fundamental} $[\partial D_{\alpha}^k]$, que es esencialmente el generador de $H_{k-1}(\partial D_{\alpha}^k)\cong \mathbb{Z}\langle [\partial D_{\alpha}^k]\rangle$, a un elemento de $H_{k-1}(X_{k-1})$. A continuación, el cociente $q: X_{k-1}\to X_{k-1}/X_{k-2} \cong \vee_{\beta}\mathbb{S}_{\beta}^{k-1}$ induce otro homomorfismo en homología 
\[
q_*: H_{k-1}(X_{k-1}) \to H_{k-1}(X_{k-1}/X_{k-2}) = H_{k-1}(\vee_{\beta}\mathbb{S}_{\beta}^{k-1}) \cong \oplus_{\beta}\mathbb{Z}\langle D^{k-1}_{\beta} \rangle = C^{CW}_{k-1}(X).
\]

A su vez, también se verifica que $\partial_{k}\circ\partial_{k+1} = 0$. Ahora que hemos visto como actuan los operadores de borde, los grupos de homología se construyen exactamente igual que en el caso simplicial.

Una vez más referimos al lector a la fuente original \cite{munoz2021} si desea aclarar algunos de los detalles técnicos más finos que no vamos a demostrar aquí. Se puede demostrar que $H_k(X)\cong H^{CW}_k(X)$ para cualquier CW-complejo $X$. Así, por transitividad se tiene que los grupos de homología de un espacio son isomorfos independientemente de la técnica que apliquemos para computarlos.

También se verifica que los grupos de homología son invariantes por homeomorfismos y, de hecho, por equivalencias homotópicas (\cite[p.~111]{hatcher2002}). Enunciamos solamente el segundo resultado en virtu de que el tipo de homotopía es también invariante por homeomorfismos.

\begin{theorem}
    Los homomorfismos $f_*:H_k(X)\to H_k(Y)$ inducidos en homología por una equivalencia homotópica $f:X\to Y$ entre dos espacios topológicos son isomorfismos.
\end{theorem}

El resultado se sigue de que dos funciones homotópicas inducen los mismos homomorfismos en homología, por tanto, una equivalencia homotópica y su inversa inducen homomorfismos inversos, luego isomorfismos.

Terminamos con unos de ejemplos de como computar la homología de un CW-complejo.

\begin{example}
    \hfill
    \begin{enumerate}[label=(\roman*)]
        \item Podemos dotar a $\mathbb{S}^n$ de una estructura de CW-complejo con una $0$-celda y una $n$-celda identificando la frontera de un $n$-disco al punto.
        
        Así, es sencillo ver que $H_0(\mathbb{S}^n)\cong \mathbb{Z}$ por ser $\mathbb{S}^n$ conexo por caminos y que $H_1(\mathbb{S}^n), \dots H_{n-1}(\mathbb{S}^n)$ son todos $\{0\}$ al ser vacíos los grupos de cadenas en esas dimensiones.

        Ahora, $H_{n}(\mathbb{S}^n)\cong \mathbb{Z}$ necesariamente ya que $Z_n(\mathbb{S}^n)=C_n(\mathbb{S}^n)\cong \mathbb{Z}$ por ser $\partial_n \equiv 0$ y $B_n(\mathbb{S}^n)=\{0\}$ por construcción del complejo de cadena.
        \item El toro $T^2$ puede ser considerado un CW-complejo con una $0$-celda $e^0$, dos $1$-celdas $e^1_1$ y $e^1_2$ y una $2$-celda $e^2$ de la siguiente manera. Tómese la construir habitual del toro plano, un cuadrado cuyos lados paralelos han sido identificados. La frontera de dicho cuadrado es la suma wedge de dos circunferencias, $\mathbb{S}^1\vee \mathbb{S}^1$, es decir, dos cricunferencias tangentes, estas serán nuestras $1$-celdas, el punto de tangencia nuestra $0$-celda. La $2$-celda será el interior del cuadrado.
        
        Ahora, $C_2(T^2)\cong \mathbb{Z}\langle e^2 \rangle \cong \mathbb{Z}$, $C_1(T^2)\cong \mathbb{Z}\langle e^1_1 \rangle \oplus \mathbb{Z}\langle e^1_2 \rangle \cong \mathbb{Z}^2$ y $C_0(T^2)\cong \mathbb{Z}\langle e^0 \rangle$.

        La $2$-celda $e^2$ se pega a las $1$-celdas a lo largo del lazo $e^1_1 e^1_2 (e^1_1)^{-1} (e^1_2)^{-1} $ donde el exponente $-1$ significa que se recorre en el sentido opuesto. En terminos de cadenas queda
        \begin{equation*}
            \partial_2(e^2) = e^1_1 + e^1_2 - e^1_1 - e^1_2 = 0.
        \end{equation*}
        Luego quedan $Z_2(T^2)=C_2(T^2)\cong \mathbb{Z}$ y $B_1(T^2)=\{0\}$. 

        Ambas $1$-celdas se pegan al $0$-esqueleto en $e^0$, ambas son lazos luego 
        \begin{equation*}
            \partial_1(e^1_1) = \partial_1(e^1_2) = e^0 - e^0 = 0.
        \end{equation*}
        Por lo que $Z_1(T^2)=C_1(T^2)\cong \mathbb{Z}^2$. De esta manera, se tiene 
        \begin{equation*}
            H_2(T^2)=Z_2/B_2 = \mathbb{Z}/\{0\} = \mathbb{Z} ~\text{ y }~ H_1(T^2)=Z_1/B1= \mathbb{Z}^2/\{0\} = \mathbb{Z}^2.
        \end{equation*}

        Y $H_0(T^2)\cong \mathbb{Z}$ por ser conexo por caminos.
        \item Por la estructura de CW-complejo dada para $\mathbb{CP}^n$ en la subsección \ref{subsec:aplicaciones} se tiene el siguiente complejo de cadena
        \begin{equation*}
            0\to C_{2n}\simeq \mathbb{Z} \to C_{2n-1}\simeq \{0\}\to \dots \to C_{2}\simeq \mathbb{Z} \to C_1\simeq \{0\}\to C_0 \simeq \mathbb{Z}\to 0 
        \end{equation*}
        donde para las dimensiones pares los grupos de cadenas son isomorfos a $\mathbb{Z}$ y para las impares a $\{0\}$. Esto significa que $Z_k = C_k$ y $B_k = \{0\}$ para todo $k$, luego los grupos de homología de $\mathbb{CP}^n$ quedan
        \begin{equation*}
            H_k(\mathbb{CP}^n) =
            \begin{cases}
                \mathbb{Z} &\text{ si } k \text{ es par},\\
                \{0\} &\text{ si } k \text{ toma otros valores}.
            \end{cases}
        \end{equation*}
    \end{enumerate}
\end{example}

\bigskip

%%%%%%%%%%%%%%%%%%%%%%%%%%%%%%%%%%%%%%%%%%%%%%%%
%%                                            %%
%%            HOMOLOGIA RELATIVA              %%
%%                                            %%
%%%%%%%%%%%%%%%%%%%%%%%%%%%%%%%%%%%%%%%%%%%%%%%%
